\section{The binary geometric-mean formula along an intertwiner}
\label{sec:geometric_mean_value_intertwine}

For positive definite matrices $A,B$ and any $p\in\R$, write
\begin{align}
  F_p(A,B)&=A^{1/2}(A^{-1/2}BA^{-1/2})^pA^{1/2}.
\end{align}
For $p\in[0,1]$ this is the weighted geometric mean. The identity
below concerns this binary formula at every real $p$.
It is auxiliary to the mean definition in Section~6
of~\cite{OpenAI2026AreaLaw}.

\begin{theorem}[Intertwining the binary mean]
  \label{thm:geometric_mean_value_intertwine}
  \lean{Matrix.geomMean_intertwine}
  \leanok
  Let $A_1,B_1$ and $A_2,B_2$ be positive definite matrices on two
  finite-dimensional complex Hilbert spaces. Suppose that a linear map
  $J$ from the second space to the first satisfies
  \begin{align}
    A_1J&=JA_2,& B_1J&=JB_2.
  \end{align}
  Then, for every $p\in\R$,
  \begin{align}
    F_p(A_1,B_1)J&=JF_p(A_2,B_2).
  \end{align}
\end{theorem}

\begin{proof}
  \leanok
  \uses{thm:real_power_intertwine}
  Theorem~\ref{thm:real_power_intertwine} gives $A_1^sJ=JA_2^s$
  for every $s\in\R$. Put $C_i=A_i^{-1/2}B_iA_i^{-1/2}$.
  These matrices are positive definite, and the two intertwining
  equations imply
  \begin{align}
    C_1J&=JC_2.
  \end{align}
  The same theorem therefore gives $C_1^pJ=JC_2^p$. Moving the two
  outer square-root factors across $J$ yields
  \begin{align}
    A_1^{1/2}C_1^pA_1^{1/2}J
      &=JA_2^{1/2}C_2^pA_2^{1/2}.
  \end{align}
  This is the asserted identity. No isometry or normalization of $J$
  is used.
\end{proof}
